% Ad Familiares 4.13 --- headnote
% ad-familiares-04-13, early August 46 BC, Rome

\textit{Cicero to P. Nigidius Figulus, written from Rome
at the beginning of August 46 BC --- the Perseus dateline
reads \textit{Scr. Romae ui. in. Sext. a. 708 (46)},
\textit{vi.\ in.\ Sext.}\ marking it as the early days of
Sextilis (August). Nigidius, a learned senator and once
praetor (58 BC), the author of a vast body of work on
grammar, theology, and natural philosophy, had taken the
Pompeian side in the civil war and was now in exile, his
estates confiscated. Caesar had not yet recalled him; he
would in fact die in exile in 45 without seeing Rome again.
Cicero himself, recently received back at Rome under
Caesar's clemency after Pharsalus, is writing from a
position of restored personal safety without restored
political freedom, and the letter is conscious throughout
of that asymmetry.}

\textit{The letter is structured as a writer's diagnosis of
his own writer's block. The opening sections work through,
and reject in turn, the available kinds of letter --- the
cheerful letter that the times have stripped away, the
promise of practical aid that Cicero in his own
diminishment cannot make, the gossiping letter about
shared friends now ruined or in flight --- before settling
on the only kind left, the consolatio, which he proceeds
to address to a man who is himself a master of the genre.
The catalogue of Nigidius's resources --- ``what your
gravity, the elevation of your spirit, the life you have
led, your studies, the arts in which you have flourished
since boyhood demand of you'' --- gives a sense of the man,
and the closing pledge to ``force my way into the company
of the man himself, which my own self-respect has hitherto
barred to me'' acknowledges that even mere access to Caesar
is now a thing Cicero has not yet brought himself to
attempt. The promised intervention bore no fruit in
Nigidius's lifetime.}
